\BeleskeID{08-s067}
% element: 08-s067:e05
% element: 08-s067:notes
U stacionarnom režimu odnos izlaza i ulaza određuje otporni razdelnik. Brzu promenu određuje kapacitivni razdelnik. Kompenzacija zahteva da se ti odnosi poklope, pa izlaz nema integratorsko zaobljavanje ivice. Jednačina odnosa se sređuje ovako:
\[\frac{C1}{C1+C2}=\frac{R2}{R1+R2}\Rightarrow C1(R1+R2)=R2(C1+C2)\Rightarrow C1R1=C2R2.\]
Pri podkompenzaciji početna promena je manja od stacionarne i izlaz prilazi konačnom nivou odozdo. Pri prekompenzaciji početna promena je veća, pa se izlaz vraća ka konačnom nivou odozgo. Dijagrami su kvalitativni; ne predstavljaju izmerene numeričke karakteristike.

Pri izjednačavanju početnog i krajnjeg izlaznog napona, zamenom odnosa deljenja dobija se
\[\frac{C1}{C1+C2}u_i(t_0^+)=\frac{R2}{R1+R2}u_i(\infty).\]
